\documentclass[11pt]{article}

% Shared preamble for the SKJS soundness corpus (main paper + deep dives).
% Loaded via \input in each document's preamble. Documents may define
% additional packages, theorem environments, and macros AFTER this input,
% but must not redefine anything here.
\usepackage{amsmath}
\usepackage{amssymb}
\usepackage{amsthm}
\usepackage{mathtools}
\usepackage{stmaryrd}
\usepackage[margin=1in]{geometry}
\usepackage{tcolorbox}
\usepackage{booktabs}
\usepackage{array}
\usepackage{longtable}
\usepackage[hyphens]{url}
\urlstyle{tt}

% Long typewriter tokens (paths, identifiers) defeat hyphenation;
% allow inter-word stretch instead of overfull lines.
\emergencystretch=2em

\newtheorem{theorem}{Theorem}
\newtheorem{lemma}[theorem]{Lemma}
\newtheorem{definition}[theorem]{Definition}
\newtheorem{corollary}[theorem]{Corollary}
\newtheorem{axiom}[theorem]{Axiom}
\newtheorem{obligation}[theorem]{Obligation}

% Summary box for key definitions
\newtcolorbox{keydef}{
  colback=gray!5,
  colframe=gray!50,
  boxrule=0.5pt,
  arc=2pt,
  left=4pt,right=4pt,top=2pt,bottom=2pt
}

% Core semantic notation shared by every document
\newcommand{\sem}[1]{\llbracket #1 \rrbracket}
\newcommand{\Val}{\mathsf{Val}}

% Implementation anchor: file / class / function at the pinned revision
\newcommand{\impl}[1]{\par\noindent{\footnotesize\textsf{impl:} \texttt{#1}}}

% Single source for the modeled revision. Every document's \date uses this;
% the ledger and tooling read the same values.
\newcommand{\skjspinhash}{9108f339}
\newcommand{\skjspindate}{2026-07-02}
\newcommand{\skjspin}{Models revision \texttt{\skjspinhash} of \texttt{main} (\skjspindate)}



\newtheorem{assumption}[theorem]{Assumption}
\newtheorem{finding}[theorem]{Finding}
\newtheorem{extension}[theorem]{Extension}

% Core semantic notation (copied from the main paper)
\newcommand{\ObjVal}{\mathsf{ObjVal}}
\newcommand{\Tags}{\mathsf{Tags}}
\newcommand{\tagfn}{\mathsf{tag}}
\newcommand{\rtag}{\mathsf{rtag}}
\newcommand{\SType}{\mathsf{SType}}
\newcommand{\STypes}{\mathsf{STypes}}
\newcommand{\TType}{\mathsf{TType}}
\newcommand{\Clauses}{\mathsf{Clauses}}
\newcommand{\Clause}{\mathsf{Clause}}
\newcommand{\TInsts}{\mathsf{TInsts}}
\newcommand{\TParam}{\mathsf{TParam}}
\newcommand{\TFields}{\mathsf{TFields}}
\newcommand{\TField}{\mathsf{TField}}
\newcommand{\Subst}{\mathsf{Subst}}
\newcommand{\SVOpen}{\mathsf{SVOpen}}
\newcommand{\SVClosed}{\mathsf{SVClosed}}
\newcommand{\clause}[2]{\langle #1 \mid #2 \rangle}


% Local notation
\newcommand{\issub}{\mathsf{isSub}}
\newcommand{\STOk}{\mathsf{STOk}}
\newcommand{\STErr}{\mathsf{STError}}
\newcommand{\vals}{\mathsf{values}}
\newcommand{\ver}{\mathsf{version}}
\newcommand{\semn}[2]{\llbracket #1 \rrbracket^{#2}}

\title{Deep Dive D --- Subtyping Validation Soundness:\\
\texttt{Clauses::isSub} and Its Strata}
\author{}
\date{\skjspin.\\
Companion to \emph{Refinement and Inference Growth in SKJS}.}

\begin{document}
\maketitle

\begin{abstract}
This deep dive addresses the largest unproved component of the main
soundness paper: Obligation \textbf{L-SUB} (\emph{Validation Soundness}),
the claim that acceptance by the subtyping validator implies semantic
inclusion under the current substitution state. Part~I proves the
architecture: the \texttt{isSub} dispatch, the two memoization caches
(\texttt{subst.seenSub} and the version-keyed \texttt{tenv.seenSub}), the
growth pre-passes, restoration after failed checks, and the coinductive
treatment of recursive types. A stratified-semantics theorem reduces the
recursive case to soundness of each individual representation layer. Part~II
establishes the layers supported by the main paper's membership semantics:
primitives, $\STypes$, DNF clause covering, object-branch matching, and the
getter, optional-field, index, and access rules of $\TFields$. Setter rules
remain conditional on \textbf{L-MUT}; function rules are checked against the
behavioral signature semantics recorded by \textbf{E-SIG}.

The analysis finds four places where acceptance is not justified by those
semantics. A callable can pass through \texttt{Function} and \texttt{\{\}}
into a type marked non-callable (\textbf{F-CABSENT}); the array brand is never
compared (\textbf{F-BRAND}); one path accepts after dropping a type-parameter
constraint; and subtyping between two implicit shape receivers, or two
implicit class receivers, does not compare their receiver types
(\textbf{F-THIS-IMPLICIT}). The registry separates these findings and the
remaining conditional cases from the parts of L-SUB established here.
\end{abstract}

\tableofcontents

%==========================================================================
\section{Purpose: the Row Addressed}
\label{sec:purpose}

The main paper's assumption ledger carries one row for the entire
validator:

\begin{keydef}
\textbf{L-SUB} (Obligation ``Validation Soundness''; the main paper's
table records the split status this paper establishes):
\texttt{Clauses::isSub} (with its $\STypes$, object-branch, field, index,
and function strata, including getter/setter variance and the
\texttt{seenSub}/\texttt{tenv.seenSub} cycle caches keyed by substitution
version) is sound in the sense of the main paper's Definition
``Subtyping Check Architecture'': if \texttt{Clauses::isSub} accepts under
state $\rho$, then $\sem{t_1}_\rho \subseteq \sem{t_2}_\rho$ for the
expanded views. \emph{``This is the largest unproved component of the
document.''}
\end{keydef}

\noindent
This paper is the discharge attempt. The outcome, in one paragraph: the
\emph{architecture} (caches, dispatch, coinduction) is proved outright,
modulo two small named assumptions; the \emph{covering skeletons}
($\STypes$, $\Clauses$, $\mathsf{TObject}$) and the \emph{read-direction}
field rules are proved against the main paper's membership semantics; the
\emph{write-direction} rules are proved conditional on the mutation model
(\textbf{L-MUT}); \emph{function} subtyping is proved against a new
behavioral signature semantics that we propose as an extension of the main
paper's object semantics (whose opaque signatures cannot even express the
claim); and two arms of the field stratum --- the lax
$(\mathsf{CUnknown},\mathsf{CAbsent})$ callability arm and the array
\emph{brand} --- are found to be unsound relative to
the main paper's stated semantics, each composing with a refinement rule
whose proof in the main paper leans on exactly
the semantic clause the subtyping arm ignores --- the first witnessed
end-to-end, the second by composite sketch
(Findings~\ref{find:cabsent} and~\ref{find:brand}).
Section~\ref{sec:status} gives the resulting ledger split.

\paragraph{Relation to the rest of the series.}
We build on three dives (exact rows in this paper's Imports table):
Deep Dive~B (\texttt{lattice-operations.tex}) for the union/intersection
results --- union proved exact, and the $\bot$-clause refuted in its
live strata (incomparable signatures and
beyond-precision num-spelling, \textbf{L-INTER-BOT});
Deep Dive~C (\texttt{type-expansion.tex}) for \textbf{L-EXP} scoped to
expansion modes \texttt{None}/\texttt{Covariant} (its theorem T-EXP),
which explicitly delegates the \emph{contravariant} expansion read to this
paper (\S\ref{sec:variance}); and Deep Dive~E
(\texttt{subst-state-discipline.tex}) for the substitution state
discipline: monotone reads and their transport theorem proved,
with the named residue MR-1 (the freeze-window memo hazard, which lives in
\texttt{subst.seenSub} and is \emph{inherited} here, \S\ref{sec:substseen}),
and the refutation of the naive ``version is a state hash'' reading of the
\texttt{tenv.seenSub} key, whose corrected argument we complete in
\S\ref{sec:tenvseen}. All code references are to revision
\texttt{\skjspinhash}; the primary sources are
\texttt{src/TTypeSubtyping.sk} (all 2203 lines audited),
\texttt{src/TTypeSubst.sk}, and the \texttt{isSub} helpers in
\texttt{src/TType.sk}.

%==========================================================================
\section{Minimal Recap}
\label{sec:recap}

\paragraph{Representation.}
From the main paper (\S3): a $\TType$ is a pair $(I, C)$ of an instance
layer $I \in \TInsts$ (type variables \texttt{TVar}, unexpanded named
instantiations \texttt{TSInst}, unexpanded intersections \texttt{TInter},
delayed array literals \texttt{TArrayLit}) and clauses $C \in \Clauses$, a
DNF: a map from var-sets to clauses $\clause{V}{S}$ with
$V \subseteq_{\mathrm{fin}} \TParam$ and $S \in \STypes \cup \{\top\}$.
$\STypes$ is tag-indexed, at most one $\SType$ per tag; $\SType$s are
units, literal-or-full primitives, and $\mathsf{object}(B)$ with branches
$b \in \TFields$ carrying named string/symbol fields, per-kind indexes,
call/new \emph{callability modes} --- the four-point lattice
$\mathsf{Callability} =
\{\mathsf{CAbsent}, \mathsf{CUnknown}, \mathsf{CFunction},
\mathsf{CPresent}(\sigma)\}$, ordered
$\mathsf{CPresent} \leq \mathsf{CFunction} \leq \mathsf{CUnknown}$ and
$\mathsf{CAbsent} \leq \mathsf{CUnknown}$
(the main paper's Definition ``Callability'';
\texttt{src/TType.sk:2217-2237}) --- and the array brand
$\mathsf{isArr} \in \{\mathsf{None},\mathsf{Some}(\mathsf{true}),
\mathsf{Some}(\mathsf{false})\}$. Two invariants are enforced at
$\TFields$ construction (\texttt{src/TType.sk:2266-2281}): an
array-branded branch lowers an unconstrained mode to $\mathsf{CAbsent}$,
and a constructible branch has $\mathsf{call} \geq \mathsf{CFunction}$
(ECMA-262: every value implementing $[[\mathrm{Construct}]]$ implements
$[[\mathrm{Call}]]$).
Semantics ($\sem{\cdot}_\rho$, indexed by the substitution state $\rho$):
\[
\sem{(I,C)}_\rho = \sem{I}_\rho \cup \sem{C}_\rho, \qquad
\sem{C}_\rho = \bigcup_{c \in C}\sem{c}_\rho, \qquad
\sem{\clause{V}{S}}_\rho = \bigcap_{v \in V}\sem{v}_\rho \cap \sem{S}_\rho,
\]
with $\sem{S} = \bigcup_{s\in S}\sem{s}$ and object membership per the
main paper's Definition ``Object-Type Semantics'': required/optional
fields constrain values via $\mathsf{get}$-types, indexes constrain the
props they \emph{cover} (string index covers both string- and numeric-kind
string keys; number index covers numeric-kind keys; symbol index covers
symbol keys), $\mathsf{call}(b)=\mathsf{CAbsent}$ entails
$\neg\mathsf{callable}(o)$, $\mathsf{CFunction}$ and
$\mathsf{CPresent}(\sigma)$ entail $\mathsf{callable}(o)$,
$\mathsf{CUnknown}$ entails nothing --- dually for $\mathsf{new}$ and
$\mathsf{constructible}$ --- and
$\mathsf{isArr}(b)=\mathsf{Some}(x)$ entails $\mathsf{isArrayV}(o)=x$.
Signatures $\sigma$ are \emph{opaque} to membership --- until
\S\ref{sec:tfun}, where we refine this.

\paragraph{Check architecture.}
The main paper's Definition ``Subtyping Check Architecture'':
\texttt{TType::isSub}$(t_1,t_2)$ runs
(1) $\mathsf{growTVars}(t_1,t_2,\mathsf{contra})$,
(2) $\mathsf{growTVars}(t_2,t_1,\mathsf{cov})$,
(3) expansion of both sides and validation by \texttt{Clauses::isSub};
on failure the substitution is restored. The Guidance-vs-Validation
theorem there shows that \emph{given} \textbf{L-SUB}, no imprecision of
growth can cause unsound acceptance: growth selects the state, validation
decides under it. This paper is about the ``validation decides'' half.

\paragraph{Target statement.}
Throughout, \emph{acceptance} means the top-level
\texttt{TType::isSub}$(\rho, t_1, t_2)$ returning $\STOk$, and the target
is:
\begin{keydef}
\textbf{L-SUB (restated).} If \texttt{TType::isSub}$(t_1,t_2)$ returns
$\STOk$, evolving the substitution $\rho \to \rho^\ast$, then
$\sem{t_1}_{\rho^\ast} \subseteq \sem{t_2}_{\rho^\ast}$.
\end{keydef}
\noindent
Two scoping remarks. (i) The statement quantifies over the state
\emph{produced} by the accepting run; persistence of the inclusion under
\emph{later} evolutions is Deep Dive~E's monotone-reads theorem, not this
paper. (ii) For types containing inference variables, $\sem{t}_\rho$
depends on the semantics assigned to a variable occurrence; we make the
required directional reads explicit in \S\ref{sec:variance} rather than
burying them in the notation.

\paragraph{Result vocabulary.}
\texttt{STResult} is $\STOk$ or an $\STErr$ tree. The error hierarchy
(several hundred of the file's 2203 lines) exists to \emph{render} failures;
it never influences acceptance. We audit it only for state effects:
\texttt{toError} renderers take the substitution and may expand types for
printing, but every render site on the acceptance path occurs strictly
after the verdict is fixed, and \texttt{originSuggestions} explicitly
works on a \texttt{subst.clone()}. Rendering is therefore out of scope for
L-SUB, with one caveat folded into the audit of
Lemma~\ref{lem:restore}.
\impl{TTypeSubtyping.sk / STError hierarchy, originSuggestions}

%==========================================================================
\part{Architecture and Caches (Proved)}

\section{The Dispatcher and Its Acceptance Points}
\label{sec:dispatch}

\begin{definition}[Dispatch]
\label{def:dispatch}
\texttt{TType::isSub}$(u_1,u_2)$ is a thin wrapper: (a) if $u_1 = u_2$,
accept; (a$'$) \emph{reflexive-\texttt{TSInst} fast path}: if $u_1$ and
$u_2$ are clause-free single \texttt{TSInst}s of the same name value and
equal arity with each argument pair equal --- syntactically, or $u_1$'s
argument a variable that \texttt{peekRaw}-resolves to $u_2$'s --- accept;
(b) if $(u_1,u_2) \in \texttt{subst.seenSub}$, accept; (c) call
\texttt{isSub\_}, and on $\STOk$ insert $(u_1,u_2)$ into
\texttt{subst.seenSub}. \texttt{isSub\_}$(u_1,u_2)$ in order:
(1) if $u_2.\mathsf{isTop}()$, accept;
(2) if $(\ver, u_1, u_2) \in \texttt{tenv.seenSub}$, accept; otherwise
    --- after the head-depth refusal of
    Definition~\ref{def:headdepth} --- record $(\ver, u_1, u_2)$ in the
    \emph{local} \texttt{tenv} copy;
(3) if $u_1$ is a single \texttt{TInter}$(\mathit{tys})$ with
    $u_2 \in \mathit{tys}$, accept;
(4) if \texttt{ArrayLitTyping::subtyping} claims the pair, return its
    verdict;
(5) save; run the two growth pre-passes; expand $u_1$ covariantly and
    $u_2$ contravariantly; validate with \texttt{Clauses::isSub};
    restore on failure.
\impl{TTypeSubtyping.sk / TType::isSub, TType::isSub\_}
\end{definition}

There are thus exactly eight ways to accept: (a), (a$'$), (b), (1)--(4)
accepting, and (5) validating. We prove (a), (1), (3) here; (a$'$) is
examined here and carries a \emph{violated} transport premise, witnessed
(Lemma~\ref{lem:fastpath}, registry row \textbf{L-SUB-TSINST}); (b) is
\S\ref{sec:substseen}; (2) in \S\ref{sec:tenvseen} and
\S\ref{sec:coinduction}; (5) is Part~II. (4) is \emph{delegated}:

\begin{obligation}[O-ARRAYLIT: Delayed Array-Literal Subtyping]
\label{obl:arraylit}
\texttt{ArrayLitTyping::subtyping} decides pairs involving delayed
array-literal forms (committing their tuple-vs-array shape) before the
general machinery runs. Its acceptance must imply inclusion, and its
failure paths must restore (the code saves/restores around its probes).
Deep Dive~C's assumption \textbf{E-ARR} covers the \emph{expansion} of
committed literals; the \emph{decision procedure} itself remains open.
Status: \textsc{open}.
\impl{TArrayLitTyping.sk / ArrayLitTyping::subtyping}
\end{obligation}

\begin{assumption}[A-HASH: Structural Identity via Hashing]
\label{as:hash}
\texttt{TType::compare} and \texttt{TType} equality are implemented as
comparison of the 64-bit \texttt{XXHash} computed by the smart
constructors (\texttt{TType.sk / compare}). Every set and map keyed by
$\TType$ in this paper --- the $u_1 = u_2$ fast path, both
\texttt{seenSub} caches, \texttt{seenExpansion} --- therefore identifies
types up to hash collision. We assume no two distinct constructed types
collide. This is the standard hash-consing wager; it is stated once here
and not repeated at each use. The reflexive-\texttt{TSInst} arm (a$'$) adds
a second identity notion --- equality of name \emph{value} plus
\texttt{peekRaw}-resolution of argument variables --- finer than raw hash
equality, analyzed directly in Lemma~\ref{lem:fastpath} rather than
resting on this wager. (A collision would be a \emph{soundness}
event: the fast path would equate distinct types.) Status:
\textsc{assumed}.
\end{assumption}

\begin{lemma}[Fast Path, Top, and the TInter Shortcut]
\label{lem:fastpath}
Acceptances (a), (1), (3) are sound. Acceptance (a$'$) satisfies the
verdict-state inclusion when the resolved argument variables are
unconstrained; its transport premise is \emph{violated}, with an
end-to-end witness (registry row \textbf{L-SUB-TSINST}).
\end{lemma}
\begin{proof}
(a) $u_1 = u_2$ (modulo A-HASH) gives $\sem{u_1}_\rho = \sem{u_2}_\rho$
for any $\rho$: the denotation is a function of the term and the state
(of the term \emph{including its position's file} --- name resolution
keys on the referencing file --- a scope shared by every
position-blind identity in this paper and folded into A-HASH's wager),
and no state evolution occurs on this path.
(a$'$) $u_1,u_2$ are clause-free with single instances
$\mathsf{TSInst}(n_1,\bar t_1)$, $\mathsf{TSInst}(n_2,\bar t_2)$; the arm
accepts only when $n_1.\mathsf{value} = n_2.\mathsf{value}$,
$|\bar t_1| = |\bar t_2|$, and for each $i$ either $t_{1,i} = t_{2,i}$ or
$t_{1,i}$ is a lone variable with
$\texttt{peekRaw}(t_{1,i}) = t_{2,i}$. Three facts bound what this
buys.
\emph{(Verdict state, unconstrained.)} \texttt{peekRaw} returns the
variable's current binding without the constraint intersection of
\texttt{Subst::get}; for an occurrence with
$\mathit{tv}.\mathsf{constraints} = \emptyset$ it denotes exactly
$\sem{t_{2,i}}_\rho$, so pointwise-equal arguments give
$\sem{u_1}_\rho = \sem{u_2}_\rho$ at the verdict state.
\emph{(Constrained occurrences.)} With
$\mathit{tv}.\mathsf{constraints} = C \neq \emptyset$ the occurrence
denotes $\sem{t_{2,i}}_\rho \cap \bigcap\sem{C}$, not
$\sem{t_{2,i}}_\rho$: the equality --- which is what lets the arm ignore
the variance of the argument position --- fails as stated, leaving only
an inclusion that invariant and contravariant argument positions do not
license.
\emph{(Transport.)} The equality is read off an \emph{unsealed}
$\SVOpen$ candidate. The slow path's covariant grow seals what it reads
(\texttt{isSub\_} step (2); \texttt{Subst::get} covariant,
$\SVOpen \to \SVClosed$, \texttt{TTypeSubst.sk:218-227}); the arm
performs the exact-value read \emph{without} the seal --- precisely the
read discipline the main paper's monotone-reads obligation forbids for
open candidates --- so a later contravariant grow widens the binding and
the verdict-state equality is not transported. This is not hypothetical:
an inferred \texttt{Box<number>} aliased through the arm, then widened
by \texttt{b.set("s")}, yields an accepted program in which a value
typed \texttt{number} holds the string \texttt{"s"} at runtime (issue
draft \path{issues/reflexive-tsinst-fastpath-unsealed-read.md}, runnable
repro). The arm returns $\STOk$ without expanding or growing
(guidance-free, no state write, not memoized) --- which is exactly why
nothing seals the peeked variable. Fix directions: seal the peeked
variables (a covariant \texttt{get}) before accepting; restrict the arm
to $\SVClosed$ bindings; or retire the arm.
(1) $u_2.\mathsf{isTop}()$ holds iff the clause map is the singleton top
clause $\clause{\emptyset}{\top}$ (\texttt{Clauses::isTop}), whose
denotation is $\Val$; note \texttt{TType::create} deletes the instance
layer when the clauses are top, so $\sem{u_2}_\rho = \Val \supseteq$
anything, and again no evolution occurs.
(3) requires $\sem{\mathsf{TInter}(\mathit{tys})}_\rho \subseteq
\sem{u_2}_\rho$ when $u_2 \in \mathit{tys}$ and $u_1$'s clause part is
empty (\texttt{getSingleInter} demands a lone instance). An unexpanded
intersection denotes at most the intersection of its conjuncts ---
this is the instance-layer counterpart of Deep Dive~C's expansion
semantics for \texttt{TInter} (its \texttt{TInter::expand} computes the
conjunction), and we record it as the interface fact
\emph{A-INTERSEM}: $\sem{\mathsf{TInter}(\mathit{tys})}_\rho \subseteq
\bigcap_{t \in \mathit{tys}} \sem{t}_\rho$. Given A-INTERSEM the shortcut
is immediate, and it performs no state write.
\end{proof}

\subsection{Restore-on-failure}
\label{sec:restore}

\begin{lemma}[Error Paths Restore]
\label{lem:restore}
Every function of the \texttt{isSub} family returns $\STErr$ only with
the substitution equal to its entry state. Audit, function by function:
\texttt{TType::isSub\_} saves before growth and restores in its single
error arm; \texttt{Clauses::isSub}, \texttt{STypes::isSub},
\texttt{TObject::isSub}, \texttt{TFields::isSub}, and \texttt{TFun::isSub}
each save at entry and restore before returning any error
(\texttt{TObject::isSub} additionally save/restores around each failed
branch pairing); \texttt{TPrimValue::isSub} takes no substitution;
the probes inside \texttt{growTVars} (\texttt{isCovered}, the
constraint-admissibility loop) save/restore around every speculative
call; \texttt{TFields::isSub\_} and \texttt{TFun::isSub\_} may leave
partial growth behind, but both are reached only through their restoring
wrappers. The one non-restoring mutation on an error path is the
\texttt{tenv.seenSub} record of step (2) --- which is a write to a
\emph{local copy} of a frozen value, not to shared state
(Lemma~\ref{lem:pathscope}) --- and the \texttt{seenExpansion} cache
fills inside \texttt{collectTVars}/\texttt{expandToVectors}, which
\emph{are} rolled back by \texttt{restore} (the cache is part of
\texttt{SubstState}) and are in any case semantics-preserving by Deep
Dive~C. This discharges, for the subtyping family, the corresponding
slice of the main paper's Obligation ``No Commitment on Failure''.
\impl{TTypeSubtyping.sk / file-head INVARIANT comment; all isSub wrappers}
\end{lemma}

\subsection{The proved-pair memo: \texttt{subst.seenSub}}
\label{sec:substseen}

\begin{definition}[Memo Discipline]
\label{def:memo}
\texttt{subst.seenSub} is a pair set inside \texttt{SubstState}. It is
written at exactly one site --- the wrapper's $\STOk$ path --- and is
captured and restored by \texttt{SavedSubst} together with
$\vals$/$\ver$/the operation caches.
\impl{TTypeSubst.sk / SubstState.seenSub; TTypeSubtyping.sk /
TType::isSub}
\end{definition}

\begin{theorem}[Memo Soundness]
\label{thm:memo}
Assume every non-memo acceptance is sound (the rest of this paper). Then
a \texttt{subst.seenSub} hit is sound, \emph{conditional on} Deep Dive~E's
monotone-reads theorem with its named residue MR-1.
\end{theorem}
\begin{proof}
Two properties combine. \emph{Timeline consistency}: because the set is
part of the saved state, an entry present at query time was inserted on
the current (never-rolled-back) timeline --- entries inserted during
exploration that failed were removed by the very \texttt{restore} that
discarded the exploration (Lemma~\ref{lem:restore}). So the entry
witnesses an actual accepting run $\rho_i \to \rho_i^\ast$ with
$\sem{u_1}_{\rho_i^\ast} \subseteq \sem{u_2}_{\rho_i^\ast}$, and every
state evolution between $\rho_i^\ast$ and the query state $\rho_q$ was a
sequence of legal transitions on the same timeline.
\emph{Transport}: Deep Dive~E's monotone-reads theorem (registry row
E-TRANSPORT) carries the
validated inclusion across those transitions. Its residue applies
verbatim here: MR-1 is precisely the case of a memo recorded while
\texttt{freezeCandidates} was set --- \texttt{Subst::get} then returns an
\emph{open candidate's} bound without sealing, so the recorded inclusion
can quote a lower bound that later grows on the $u_1$ side. We do not
re-derive it; the memo inherits MR-1 as a conditional hypothesis. Note
the contrast with \texttt{tenv.seenSub}: no version key is needed here
exactly because failed timelines are rolled back and successful growth
is what the transport theorem is about.
\end{proof}

\paragraph{The freeze-window residue.} The conditional hypothesis fails
in the implementation: MR-1 is \textsc{violated}, witnessed end-to-end
by an order-dependent unsound acceptance through a stale
\texttt{seenSub} entry recorded in a freeze window (deep dive E; issue
draft \path{issues/freeze-window-stale-memo-caches.md}). The lemma
above therefore holds only \emph{outside} freeze windows:
\texttt{subst.seenSub} is unsound inside them until the memo is
version-keyed or freeze-window writes are suppressed. The rest of Part~I
is unaffected — the proof isolates the failure to exactly this
hypothesis.

\subsection{Growth pre-passes are guidance}
\label{sec:growth}

Step (5) of Definition~\ref{def:dispatch} runs
$\mathsf{growTVars}(u_1,u_2,\mathsf{contra})$ then
$\mathsf{growTVars}(u_2,u_1,\mathsf{cov})$ before validating. By the main
paper's Guidance-vs-Validation theorem, these passes require no soundness
proof of their own: they only choose the state $\rho'$ under which
\texttt{Clauses::isSub} decides, and the failure path restores
(Lemma~\ref{lem:restore}). Three details of the audit are still worth
recording, because each is a place where a reader might suspect a hidden
acceptance. (i) The \texttt{isCovered} probes and the
constraint-admissibility checks inside \texttt{growTVars} call
\texttt{TType::isSub} speculatively; their verdicts influence only
\emph{which} disjuncts are grown, never the outer verdict, and each probe
restores. (ii) The admissibility loop runs under
\texttt{freezeCandidates}; by the main paper's never-widen lemma this
makes it a query on already-determined candidates --- but its
\emph{successful} probes may insert \texttt{subst.seenSub} memos while
frozen, which is one of the two entry points of the inherited MR-1
hazard (Theorem~\ref{thm:memo}). (iii) The
\texttt{constraintErr} value threaded out of growth is used only to decorate
the error message with hints on the failure path. (iv) When both sides
are tuples of the same arity, the structural descent seeds each element
variable from its positional pair only, not from the position-unioning
method members (\texttt{pop}, \texttt{at}, \ldots) --- a precision
choice inside guidance: it changes which state validation decides under,
never the verdict logic, so the lemmas of Part~II are unaffected.
\impl{TTypeSubtyping.sk / growTVars, collectTVars, isSub\_ (tuple
positional seeding)}

\section{The Cycle Cache: \texttt{tenv.seenSub} and Version Keying}
\label{sec:tenvseen}

The interesting cache is the second one. Its entry is recorded \emph{at
entry}, before the verdict is known --- the implementation comment says
it plainly: ``if a cycle occurs, the subtyping relation is considered
correct (by co-induction).'' Soundness therefore needs two arguments: a
\emph{state} argument --- a cache hit must not smuggle in a judgement
recorded under different bindings --- proved in this section, and a
\emph{logical} argument --- why admitting a pair before it is proved is
sound at all --- which is \S\ref{sec:coinduction}.

\subsection{The write audit: every binding write bumps the version}

\begin{lemma}[Version Audit]
\label{lem:version}
In \texttt{TTypeSubst.sk}, every assignment to \texttt{Subst.values} is
immediately preceded by \texttt{version += 1}, and these are the only
version bumps besides \texttt{gen}/\texttt{genForField}/%
\texttt{genForLiteralCommitment} (which bump without writing
\texttt{values}) and \texttt{restore} (which sets $\ver$ and $\vals$
jointly from a snapshot). Consequently, between two moments with no
intervening \texttt{restore}, $\ver$ unchanged implies $\vals$ unchanged.
\end{lemma}
\begin{proof}
By inspection of the four write sites: \texttt{grow} bumps at its single
write (and its idempotent arm --- $t_2 \vee t = t_2$ --- returns
\emph{before} both the bump and the write, a fact
\S\ref{sec:coinduction} uses for cycle convergence); \texttt{get} bumps at
both sealing writes (fresh-to-top and open-to-closed); no other method
writes \texttt{values}. \texttt{TVar::setConstraint} is \emph{not} a
binding write: as established in Deep Dive~E, it is a functional update
producing a new \texttt{TVar} occurrence (and the new occurrence hashes
differently, so cache keys distinguish it --- A-HASH). Callers outside the
file (\texttt{TArrayLitTyping.sk}, \texttt{TVariadicTyping.sk},
\texttt{TypeCheckUtils.sk}) reach \texttt{values} only through
\texttt{grow}/\texttt{get}, per the main paper's write-site audit.
\end{proof}

\subsection{Why the naive reading fails, and what replaces it}
\label{sec:naive}

The comforting reading --- ``the key $(\ver, u_1, u_2)$ identifies the
substitution state, so a hit replays a judgement made under the same
bindings'' --- is \emph{false}, as Deep Dive~E established: \texttt{restore}
rewinds the counter, so after \texttt{save} at $v$, a write to $v{+}1$,
a restore to $v$, and a \emph{different} write to $v{+}1$, two distinct
binding maps have carried the same version number. Version numbers are
reused across divergent exploration branches (and \texttt{gen} even reuses
\texttt{TVar} identifiers across them). If \texttt{tenv.seenSub} were a
global mutable table, entries recorded on a discarded branch would remain
reachable on the surviving branch at colliding keys, and --- worse ---
entries recorded at entry by checks that then \emph{failed and restored}
would be hit again at the identical rewound state, converting a rejection
into a silent acceptance. Our first proof attempt followed the naive
reading and foundered on exactly these two scenarios; we record that as a
methodological finding: \emph{the version key alone proves nothing}. What
actually makes the cache sound is a pair of structural facts about how
\texttt{tenv} and the saved substitutions travel, proved next.

\subsection{Path scoping}

\begin{lemma}[Path Scoping]
\label{lem:pathscope}
At any moment during a top-level \texttt{TType::isSub} run, the
\texttt{tenv.seenSub} value visible to the currently executing
\texttt{isSub\_} invocation is exactly the ambient set of the top-level
entry (empty in practice) plus one record
$(\ver_A, a_1, a_2)$ for each \emph{strict ancestor} \texttt{isSub\_}
invocation $A$ on the current call stack, recorded at $A$'s entry.
Records made by invocations that have already returned --- successfully
or not --- are not visible.
\end{lemma}
\begin{proof}
\texttt{TEnv} is a frozen class and \texttt{tenv} is passed by value
through every signature in the file. The record statement
\texttt{!tenv.seenSub = tenv.seenSub.set(...)} is Skip's functional
update of the \emph{local} binding \texttt{tenv} (contrast
\texttt{subst.!seenSub = ...}, a true field mutation on the mutable
\texttt{Subst}); the updated copy flows only into the callees of the
recording invocation --- \texttt{growTVars}, \texttt{expandToClauses},
\texttt{Clauses::isSub} and its strata --- and evaporates when that
invocation returns. A grep-complete audit shows \texttt{isSub\_} is the
only writer of the field in the codebase. Hence visibility is precisely
``my ancestors' records'': each ancestor recorded into its local copy
before descending, and no other write exists.
\end{proof}

\subsection{Record-point coherence}

\begin{lemma}[Record-Point Coherence]
\label{lem:coherence}
Suppose an \texttt{isSub\_} invocation hits the cache:
$(\ver_A, u_1, u_2) \in \texttt{tenv.seenSub}$ with $\ver = \ver_A$ now.
Then $\vals$ now equals $\vals$ at the moment ancestor $A$ recorded the
entry.
\end{lemma}
\begin{proof}
Call $A$'s record moment $t_A$, with state $(\ver_A, X_A)$. Two
sub-audits establish that every state on the current path between $t_A$
and now with version $\ver_A$ has $\vals = X_A$.

\emph{(i) Save/restore locality.} Every \texttt{SavedSubst} in the
subtyping family is a local variable, created and consumed in the same
invocation (Lemma~\ref{lem:restore}); none is passed down or stored.
Therefore every \texttt{restore} executed while $A$ is pending targets a
snapshot taken \emph{after} $t_A$ (the restoring invocation was entered
after $A$, and its save postdates its entry). The wholesale
snapshot/restore pairs outside the family
(\texttt{saveAll}/\texttt{mcreateFromSaved} at the SKStore boundary,
\texttt{clone} used by rendering and refinement) never fire inside a
pending \texttt{isSub\_}: \texttt{clone} creates a detached substitution
whose mutations do not touch the current one.

\emph{(ii) Version floor and determinism.} By induction over the
execution steps from $t_A$: every reachable state while $A$ is pending
has $\ver \geq \ver_A$, and if $\ver = \ver_A$ then $\vals = X_A$.
Writes only increase $\ver$ (Lemma~\ref{lem:version}). A restore targets
a post-$t_A$ snapshot $(v_s, X_s)$; by the induction hypothesis applied
at snapshot time, $v_s \geq \ver_A$, and if $v_s = \ver_A$ then
$X_s = X_A$; restoring re-establishes exactly that pair. \texttt{gen}
bumps without writing, preserving both claims. So at the hit,
$\ver = \ver_A$ forces $\vals = X_A$.
\end{proof}

\noindent
The two lemmas together repair the naive reading: cross-branch collisions
cannot reach the cache (discarded branches' records were never visible,
Lemma~\ref{lem:pathscope}), failed checks leave no landmine (their records
die with their stack frame), and a genuine hit is a \emph{pending
ancestor} re-encountered at \emph{bit-identical bindings}
(Lemma~\ref{lem:coherence}). What the version key buys is the second
half: without it, a descendant could hit an ancestor's record after
intervening growth, treating a judgement-in-progress under $\rho$ as
applicable under a strictly grown $\rho'$ --- the stale-cache regression
class (\#446) the key was introduced to kill. What remains to justify is
the logic of accepting a pending pair at all.
\impl{TTypeSubtyping.sk / isSub\_ (lines 1852, 1877); TEnv.sk / seenSub
comment}

\subsection{The head-depth refusal}

\begin{definition}[Head-Depth Refusal]
\label{def:headdepth}
Between the cycle-cache check and its record, \texttt{isSub\_} counts,
per ordered pair of \emph{head constructor names} (both sides lone
\texttt{TSInst}s), how many times the current path has entered that
pair, in a map \texttt{tenv.subtypingHeadDepth} carried exactly like
\texttt{tenv.seenSub} (a functional update of the local frozen-\texttt{TEnv}
copy, so Lemma~\ref{lem:pathscope}'s path-scoping applies verbatim).
Past the bound (\texttt{kMaxSubtypingHeadDepth} $= 12$) the check is
\emph{refused} with TS2589 rather than assumed related: two recursive
aliases whose arguments grow at every level never repeat a
\texttt{seenSub} entry, so without the bound their descent would be
limited by nothing.
\impl{TTypeSubtyping.sk / isSub\_ (lines 1852-1877); kMaxSubtypingHeadDepth
(line 18); TEnv.sk / subtypingHeadDepth}
\end{definition}

\noindent
For \textbf{L-SUB} the refusal is the sound direction --- it can only
reject --- and it does not perturb the coinductive argument: a
\texttt{tenv.seenSub} hit still names a pending ancestor at identical
bindings (the record still precedes the descent), and the bound merely
caps how many same-constructor entries a path accumulates before
refusing. Its cost is completeness (deep finite nestings past the bound
are rejected; the constant makes this rare), and its benefit is the
termination remark of the honesty box below: the implementation carries
its own divergence guard for the one recursion pattern the cycle cache
cannot close.

\section{Coinduction over the Value Semantics}
\label{sec:coinduction}

\subsection{What the semantics of a recursive type even is}

The main paper's object semantics is a system of \emph{equations}:
$o \in \sem{b}$ is defined in terms of $\mathsf{get}(o,p) \in
\sem{\mathsf{get}(f)}$ for field types that, for a recursive declaration
(\texttt{type L = \{ next: L \}} via \texttt{TSInst} unfolding), mention
$\sem{b}$ again. The
equations are monotone --- every recursive occurrence of a denotation in
the membership conditions is positive: field and index constraints use
$\sem{\cdot}$ on the membership side, the call/new and brand clauses are
atomic predicates on $o$, and nothing negates a membership --- so by
Knaster--Tarski both a least and a greatest solution exist, and the main
paper's definition takes the greatest one (registry row SEM-GFP, which
this paper consumes).

\begin{definition}[SEM-GFP: Greatest-Fixpoint Reading (main-paper recap)]
\label{as:gfp}
The denotation $\sem{\cdot}$ is the \emph{greatest} solution of
the membership equations. This is the reading under which cyclic values
inhabit recursive types: \texttt{const x = \{\}; x.next = x} builds a
cyclic value that TypeScript and SKJS deliberately type as \texttt{L},
and the greatest solution contains exactly the values all of whose
finite explorations are consistent with the type.
\end{definition}

Concretely we use the standard stratification. Define
$\semn{t}{0} = \Val$ and let $\semn{t}{n+1}$ be the main paper's
membership conditions with every \emph{structural} recursive membership
--- field values, index values, and (after \S\ref{sec:tfun}) signature
argument/result memberships --- evaluated at level $n$, while tag tests,
literal tests, callability, and brands are evaluated exactly. Unions,
intersections, clause var-sets, and instance layers distribute through
levels pointwise. Then $\sem{t} = \bigcap_n \semn{t}{n}$ is the greatest
solution (the defining operator preserves the intersection of this
decreasing chain because each value's membership at level $n{+}1$
depends on finitely-branching conditions at level $n$; the union over
object branches and clauses is over finite sets, so no
non-$\omega$-continuity arises on the relevant chain).

\subsection{One-step soundness, and the reduction}

\begin{definition}[One-Step Soundness of a Stratum]
\label{def:onestep}
Let $R$ be a set of $\TType$ pairs (``the accepted pairs''). A stratum
check is \emph{one-step sound} if, whenever it accepts a pair
$(x_1, x_2)$ at state $\rho$ with all of its recursive
\texttt{TType::isSub} sub-checks resolving to pairs in $R$, then for
every $n$:
\[
\big(\forall (a,b) \in R.\ \sem{a}_\rho \subseteq \semn{b}{n}_\rho\big)
\implies
\sem{x_1}_\rho \subseteq \semn{x_2}{n+1}_\rho .
\]
Part~II proves each stratum one-step sound (or names the failure). For
readability Part~II states the lemmas in the naive form ``acceptance
implies inclusion, given inclusion at the recursive positions''; each
proof visibly uses the recursive hypotheses only inside one layer of
membership conditions, which is exactly the shape
Definition~\ref{def:onestep} requires, so the stratified reading is
obtained by re-indexing, and we do not repeat the bookkeeping.
\end{definition}

\begin{theorem}[Coinductive Reduction]
\label{thm:coind}
Assume A-HASH, SEM-GFP, the state discipline of
\S\S\ref{sec:dispatch}--\ref{sec:tenvseen}, and one-step soundness of
every stratum reachable from \texttt{Clauses::isSub}. Then \textbf{L-SUB}
holds: a top-level accepting run $\rho \to \rho^\ast$ satisfies
$\sem{t_1}_{\rho^\ast} \subseteq \sem{t_2}_{\rho^\ast}$.
\end{theorem}
\begin{proof}
Consider the finite call tree of the accepting run and let $R$ be the set
of pairs of all \texttt{TType::isSub} nodes that returned $\STOk$,
\emph{each taken with the state at its verdict}. First flatten the state
indexing: every node's verdict state precedes $\rho^\ast$ on the same
never-restored timeline (failed excursions restored,
Lemma~\ref{lem:restore}), so by Deep Dive~E's monotone-reads transport
(with its MR-1 proviso, and modulo the unsealed exact-value read of arm
(a$'$), registry row \textbf{L-SUB-TSINST} --- the transport premise
that arm violates) it suffices to prove each node's inclusion at its
own verdict state; we henceforth suppress $\rho$.
Now show $\forall n.\ \forall (a,b) \in R.\ \sem{a} \subseteq \semn{b}{n}$
by induction on $n$. Base: $\semn{b}{0} = \Val$. Step: fix
$(a,b) \in R$ and $v \in \sem{a}$. The node for $(a,b)$ accepted through
one of the eight acceptance points. Fast path/top/TInter: by
Lemma~\ref{lem:fastpath} the full inclusion $\sem{a} \subseteq \sem{b}$
holds outright, and $\sem{b} \subseteq \semn{b}{n+1}$. The
reflexive-\texttt{TSInst} arm (a$'$): the verdict-state inclusion needs
unconstrained argument variables, and the flattening step above already
carries its violated transport premise, so this case --- hence the
theorem --- is conditional on registry row \textbf{L-SUB-TSINST},
exactly as the transport step is on MR-1.
Memo (\texttt{subst.seenSub}): by Theorem~\ref{thm:memo} the entry
witnesses an earlier accepting node, i.e.\ another element of $R$ with
the same pair; use its case.
Cycle cache (\texttt{tenv.seenSub}): by
Lemmas~\ref{lem:pathscope}/\ref{lem:coherence} the hit names a
\emph{pending ancestor} node with the same pair and bindings; that
ancestor either eventually accepted --- it did, since the whole run
accepted and an ancestor of an executed node has its verdict on the
accepting path --- so the pair is in $R$, and \emph{the induction is over
$n$, not over the tree}: circular use of the pair at level $n$ is
precisely what the induction hypothesis licenses.
Delegated point (4): conditional on Obligation~\ref{obl:arraylit}.
Validation (5): \texttt{Clauses::isSub} accepted with its recursive
sub-checks all $\STOk$, i.e.\ with pairs in $R$; one-step soundness of
the strata gives $v \in \semn{b}{n+1}$ from the induction hypothesis.
Finally $\sem{a} \subseteq \bigcap_n \semn{b}{n} = \sem{b}$ by SEM-GFP,
and instantiating at the root gives L-SUB.
\end{proof}

\noindent
\emph{Honesty box.} Theorem~\ref{thm:coind} is genuinely conditional on:
MR-1 (inherited
via the memo and via transport); O-ARRAYLIT; A-HASH; and one-step
soundness of every stratum --- Part~II shows the last is where the
substance and the violations live. Termination (that the cycle cache is
\emph{hit} rather than the growth diverging --- convergence relies on the
idempotent-grow early return noted in Lemma~\ref{lem:version}, and the
argument-growing recursion the cache cannot close is refused by the
head-depth bound, Definition~\ref{def:headdepth}) is a
liveness property, out of scope for L-SUB.

%==========================================================================
\part{The Strata}

\section{Primitives: Units and \texttt{TPrimValue::isSub}}
\label{sec:prim}

\begin{definition}[Primitive Checks]
Both checks are invoked only on same-tag pairs (an \texttt{invariant}
guards this; the caller guarantees it, \S\ref{sec:stypes}). Units:
accepted unconditionally. \texttt{TPrimValue::isSub}$(p_1, p_2)$ on
literal sets: $(-, \top)$ accept; $(\top, L_2)$ reject; $(L_1, L_2)$
accept iff every stored spelling in $L_1$ occurs in $L_2$ (string
equality of normalized lexemes, \texttt{lits2.contains(elt.val)}).
It touches no state.
\impl{TTypeSubtyping.sk / TPrimValue::isSub, STypes::isSub (unit arm)}
\end{definition}

\begin{lemma}[Primitive Strata are Sound --- Unconditionally]
\label{lem:prim}
(i) For same-tag units $u_1, u_2$: $\sem{u_1} = \sem{u_2}$.
(ii) $\texttt{TPrimValue::isSub}(p_1,p_2) = \STOk \implies
\sem{p_1} \subseteq \sem{p_2}$.
\end{lemma}
\begin{proof}
(i) The three unit constructors have three distinct tags
(\textsf{undef}, \textsf{void}, \textsf{null}), so same tag forces the
same unit and singleton denotations coincide.
(ii) $(-,\top)$: by the main paper's A0 every $v \in \sem{p_1}$ has
$p_1$'s tag, and $\sem{\mathsf{prim}(k,\top)}$ is the full tag set.
$(L_1, L_2)$: let $v \in \sem{p_1}$, so $v \approx
(\tagfn(k), \delta_k(\ell))$ for some $\ell \in L_1$; the spelling
$\ell$ occurs in $L_2$, hence $v \in \sem{p_2}$ by the same closure
clause. This uses only that $\delta_k$ is a \emph{function} of the
spelling --- \textbf{L-LIT}(i), the proved part --- not injectivity: a
duplicate value under two spellings can only cause spurious
\emph{rejection} (a completeness matter, L-LIT territory), never
acceptance. No
recursive positions exist, so one-step soundness is the plain statement.
\end{proof}

This stratum is the model case for the paper's staging: fully proved,
no conditions, no state.

\section{The Union Spine: \texttt{STypes::isSub}}
\label{sec:stypes}

\begin{definition}[Tag-Aligned Covering]
\texttt{STypes::isSub}$(S_1, S_2)$ runs \texttt{leq} over the tag-keyed
maps: for each member $s_1 \in S_1$, if $S_2$ has no member at
$\tagfn(s_1)$, fail (\texttt{STTooManySType}); otherwise dispatch the
same-tag pair to the unit/\texttt{TPrimValue}/\texttt{TObject} stratum.
Members of $S_2$ at tags absent from $S_1$ are ignored. Saves at entry,
restores on failure.
\impl{TTypeSubtyping.sk / STypes::isSub;
InternalKeyedArrayMap.sk / leq, iterDiffInterDiff}
\end{definition}

\begin{lemma}[$\STypes$ Covering is Sound]
\label{lem:stypes}
If \texttt{STypes::isSub}$(S_1,S_2) = \STOk$ and each accepted same-tag
pair $(s_1, s_2)$ satisfies $\sem{s_1} \subseteq \sem{s_2}$ (the
sub-stratum obligation, one-step for objects), then
$\sem{S_1} \subseteq \sem{S_2}$.
\end{lemma}
\begin{proof}
Let $v \in \sem{S_1} = \bigcup_{s \in S_1}\sem{s}$, say $v \in \sem{s_1}$.
By A0, $\rtag(v) = \tagfn(s_1)$. \texttt{leq} succeeded, so $S_2$ has a
member $s_2$ at that tag --- the keyed-map structure (at most one member
per tag, iterated by sorted key) is what makes ``the'' member
well-defined and the \texttt{invariant} in the code unreachable-else ---
and the sub-stratum accepted $(s_1, s_2)$, giving
$v \in \sem{s_2} \subseteq \sem{S_2}$. Ignoring $S_2$'s extra tags is
sound: they only enlarge the right side. One-step form: the object arm's
hypothesis is used at level $n$, delivering $\semn{s_2}{n+1} \ni v$,
and unit/primitive arms are exact at every level.
\end{proof}

Completeness is not claimed and indeed fails (a literal member of $S_1$
covered by \emph{no single} same-tag member cannot borrow from other
tags --- correct, by A1 there is nothing to borrow), but incompleteness
is not our concern. Note the check inherits \emph{no} condition from
Deep Dive~B's L-INTER-BOT refutation: no intersection is consumed on this
path.

\section{DNF Covering: \texttt{Clauses::isSub}}
\label{sec:clauses}

\begin{definition}[Clause Covering]
\texttt{Clauses::isSub}$(C_1, C_2)$: for every clause
$\clause{V_1}{S_1} \in C_1$, search $C_2$ for a clause
$\clause{V_2}{S_2}$ with (gate) $V_2 \subseteq V_1$ --- the code's
\texttt{inter1.vars.contains(inter2.vars)}, a superset test by the
\texttt{ArraySet::contains} definition --- and (body) one of:
$S_1 = S_2 = \top$: match; $S_1 \neq \top, S_2 = \top$: match;
$S_1 = \top, S_2 \neq \top$: skip this candidate;
both $\neq \top$: match iff \texttt{STypes::isSub}$(S_1,S_2)$ accepts.
Accept iff every $C_1$-clause finds a match; restore on failure.
\impl{TTypeSubtyping.sk / Clauses::isSub}
\end{definition}

\begin{lemma}[Clause Covering is Sound, Parametrically in $\TParam$]
\label{lem:clauses}
If \texttt{Clauses::isSub} accepts (with the $\STypes$ obligations of the
matched pairs), then for \emph{every} interpretation
$\iota : \TParam \to \mathcal{P}(\Val)$ of the type parameters,
$\sem{C_1}^\iota \subseteq \sem{C_2}^\iota$ where
$\sem{\clause{V}{S}}^\iota = \bigcap_{p \in V}\iota(p) \cap \sem{S}$.
\end{lemma}
\begin{proof}
Fix $\iota$ and $v \in \sem{\clause{V_1}{S_1}}^\iota$ for some
$C_1$-clause, and let $\clause{V_2}{S_2}$ be its match.
Gate: $V_2 \subseteq V_1$ gives
$\bigcap_{V_1}\iota \subseteq \bigcap_{V_2}\iota$ --- for \emph{any}
$\iota$, which is the parametricity: the algorithm never inspects a
parameter's meaning, so its verdict is uniform in $\iota$, and
instantiating parameters later (generic call sites, the
\texttt{RenameClauses} step of \S\ref{sec:tfun}) preserves the inclusion.
Body: if $S_2 = \top$, $\sem{S_2} = \Val$ and we are done. Both concrete:
Lemma~\ref{lem:stypes} gives $\sem{S_1} \subseteq \sem{S_2}$
(one-step: at level $n{+}1$ via the hypotheses at $n$). The skipped case
($S_1 = \top$, $S_2$ concrete) rejects nothing by itself; it merely
forces another candidate, and if none exists the check fails --- sound.
Union over $C_1$'s clauses completes it.
\end{proof}

Two deliberate incompletenesses, noted for contrast with soundness: a
$C_1$-clause must be covered by a \emph{single} $C_2$-clause (no covering
by unions of clauses --- e.g.\ $T \mid \texttt{string}$ on the right
never combines its parts), and the gate demands syntactic var-set
containment rather than any semantic entailment between parameters.
Both only reject.

\section{Object Branches: \texttt{TObject::isSub}}
\label{sec:tobject}

\begin{definition}[Branch Matching]
\texttt{TObject::isSub}$(o_1, o_2)$: every branch
$b_1 \in \mathsf{branches}(o_1)$ must satisfy
\texttt{TFields::isSub}$(b_1, b_2)$ for \emph{some}
$b_2 \in \mathsf{branches}(o_2)$; each failed pairing restores before the
next candidate is tried, and overall failure restores to entry.
\impl{TTypeSubtyping.sk / TObject::isSub}
\end{definition}

\begin{lemma}[Branch Matching is Sound]
\label{lem:tobject}
If \texttt{TObject::isSub} accepts and each accepted pairing satisfies
$\sem{b_1} \subseteq \sem{b_2}$ (one-step, the $\TFields$ obligation),
then $\sem{o_1} \subseteq \sem{o_2}$.
\end{lemma}
\begin{proof}
Membership in an object type is existential over branches
($o \in \sem{\mathsf{object}(B)}$ iff $o \in \sem{b}$ for some
$b \in B$), so the $\forall\exists$ shape of the algorithm matches the
semantics exactly: $v \in \sem{b_1}$ for some $b_1$, its matched $b_2$
gives $v \in \sem{b_2} \subseteq \sem{o_2}$. The per-pairing
save/restore matters for a different reason than soundness of the
verdict: a failed $b_2$-candidate may have committed growth (its field
checks grow inference variables); restoring keeps the surviving
pairing's state from being polluted by the failed one --- without it the
\emph{accepted} state $\rho^\ast$ would contain bindings from abandoned
pairings, and while Guidance-vs-Validation says that could not falsify
the verdict, it would degrade the solutions chosen. As with clauses, the
single-branch covering is incomplete (no splitting a sub branch across
super branches) but sound.
\end{proof}

\section{Fields: \texttt{TFields::isSub}}
\label{sec:tfields}

This is the widest stratum. \texttt{isSub\_} checks, in order: named
string fields; the string index pair; the \emph{effective} number index
pair; the symbol index pair; named symbol fields (with the symbol index
as their covering index); the call mode; the new mode. Named-field
pairing uses \texttt{iterDiffInterDiff} over the name-keyed maps:
$\mathsf{Left}(f_1)$ = sub-only field, $\mathsf{Right}(f_2)$ = super-only
field, $\mathsf{Both}(f_1,f_2)$ = shared name. Throughout this section
$o \in \sem{b_1}$ is fixed and we must establish each conjunct of
$o \in \sem{b_2}$ per the main paper's object semantics; recursive
\texttt{TType::isSub} calls are the one-step positions.
\impl{TTypeSubtyping.sk / TFields::isSub\_, isSub\_NamedFields}

\subsection{Shared fields ($\mathsf{Both}$)}

The code: reject if $f_1$ optional and $f_2$ required; reject if
$\neg f_1.\mathsf{access}.\mathsf{isCompatibleWith}(f_2.\mathsf{access})$;
reject if $f_1.\mathsf{set} = \bot$ (readonly) and
$f_2.\mathsf{set} \neq \bot$ (skipping the recursive check to avoid a
TVar in $f_2.\mathsf{set}$ unifying its way to spurious success --- a
nice example of growth-avoidance \emph{for precision}, not soundness);
otherwise require $f_2.\mathsf{set} \leq f_1.\mathsf{set}$
(contravariant) and $f_1.\mathsf{get} \leq f_2.\mathsf{get}$ (covariant,
against \emph{bare} $f_2.\mathsf{get}$ even when both are optional ---
exact-optional semantics).

\begin{lemma}[Shared-Field Soundness, Read Direction]
\label{lem:both-get}
The optionality gate, the access gate, and the covariant get check
suffice for the \emph{field clauses} of $o \in \sem{b_2}$ at the shared
names, given the recursive inclusions.
\end{lemma}
\begin{proof}
If $f_2$ is required, the gate makes $f_1$ required, so
$\mathsf{has}(o, \mathsf{name})$ holds and (unless phantom)
$\mathsf{get}(o,\mathsf{name}) \in \sem{f_1.\mathsf{get}} \subseteq
\sem{f_2.\mathsf{get}}$ by the accepted recursive pair. If $f_2$ is
optional, its clause constrains only the case
$\mathsf{has}(o,\mathsf{name})$; when the key is present, membership in
$b_1$ gives the value in $\sem{f_1.\mathsf{get}}$ \emph{whether $f_1$ is
optional or required} --- this is exactly why checking against bare
$f_2.\mathsf{get}$ (not $f_2.\mathsf{get} \cup \texttt{undefined}$) is
the right reading: the optional read's \texttt{| undefined} accounts for
\emph{absence}, and absence is vacuous here. The exactness comment in
the code records the same argument. Phantom-access $f_1$: its get is
$\top$, so the recursive check fails against any concrete
$f_2.\mathsf{get}$ and the pair is rejected --- restrict-only; the access
gate independently rejects phantom-vs-declared
(\texttt{FAPhantom} is compatible only with \texttt{FAPhantom}/%
\texttt{FANever}), and a phantom-vs-phantom acceptance constrains only
presence on both sides. Private fields: compatibility demands the same
owner name and file --- nominal, strictly finer than the public case, so
acceptance implies the same get/set checks ran; \texttt{FANever} as
target is vacuously safe (nothing may access it; as a membership
constraint it is an ordinary field and the get/set checks still ran).
\impl{TType.sk / FieldAccess::isCompatibleWith}
\end{proof}

\begin{lemma}[Shared-Field Soundness, Write Direction --- Conditional
(registry row L-SUB-FIELDS-SET)]
\label{lem:both-set}
The contravariant set check
($f_2.\mathsf{set} \leq f_1.\mathsf{set}$) and the readonly gate are
\emph{not} consequences of membership: the main paper's semantics
constrains values via get-types only, so
Lemma~\ref{lem:both-get} already delivers $o \in \sem{b_2}$ without
them. What they buy is the mutation model: under \textbf{L-MUT}, a write
through the $b_2$ view is checked against $f_2.\mathsf{set}$, and the
contravariant condition guarantees any such accepted value is also
legal for the underlying $f_1.\mathsf{set}$, i.e.\ writes through the
super view cannot break invariants the sub view's clients rely on
(the readonly gate is the $\bot$ instance: no value passes
$\leq \bot$, so a readonly field never acquires writers by subsumption).
Status: conditional on \textbf{L-MUT} --- if the mutation model is ever
made concrete, these checks are its proof obligations, not membership's.
\qed
\end{lemma}

\subsection{Sub-only fields ($\mathsf{Left}$) and super-only fields
($\mathsf{Right}$)}

\begin{lemma}[Width and Optional-Absence Soundness]
\label{lem:width}
The $\mathsf{Left}$ and $\mathsf{Right}$ arms are sound for membership,
given the recursive inclusions (read direction; their set-direction
checks are conditional exactly as in Lemma~\ref{lem:both-set}).
\end{lemma}
\begin{proof}
$\mathsf{Left}(f_1)$: a field only the sub declares. The super's
membership conditions mention the name only through its \emph{indexes}:
if $b_2$ has an index covering the name
(\texttt{indexForFieldName}: numeric-spelling names consult the number
index first, then the string index --- matching the covering semantics,
under which a numeric-kind key must satisfy \emph{both} when both exist;
registry row \texttt{L-SUB-IDXNAME}, \textsc{open}: the code checks
against only the first-consulted index
(\texttt{src/TType.sk:2135-2141}; call site
\texttt{src/TTypeSubtyping.sk:1172}), so a numeric-named sub field is
under-checked when the super has both a number and a string index. The
routes to the under-check are narrow, three of them verified against
the binary at the pin: a named-field-only source is rejected outright
against an index-declaring super --- the
$(\mathsf{None},\mathsf{Some})$ index-pair arm, which the
implementation's own comment records as the guard
(\texttt{src/TType.sk:2127-2134}); a declaration whose number and
string indexes carry incompatible value types is rejected at the
declaration itself; and reads of numeric-named keys resolve through the
number index --- the same index the width check consulted --- while
string-variable reads stay conservative. A validator-completeness gap,
not a witnessed unsoundness; the ``suffices''
below therefore holds only under those guards, not in general), the code
requires
$f_1.\mathsf{get} \leq \mathit{idx}_2.\mathsf{get}$: then when
$\mathsf{has}(o, \mathsf{name})$, the value lies in
$\sem{f_1.\mathsf{get}} \subseteq \sem{\mathit{idx}_2.\mathsf{get}}$,
which is what the index clause of $o \in \sem{b_2}$ demands of covered,
non-named props. If $b_2$ has no covering index, the name is
unconstrained by $b_2$ --- the semantics is extensible --- and accepting
silently is exact. (The footnote in \texttt{TIndexesBase} about
\texttt{"Infinity"}/\texttt{"NaN"} spellings classifying as string-kind
is a key-\emph{kind} discrepancy against TypeScript, argued unobservable
there; for us it only changes \emph{which} index the check runs against,
and both candidate indexes' clauses are enforced by the index-pair
checks, so soundness is insensitive to the classification.)
$\mathsf{Right}(f_2)$ required: rejected --- necessary, since extensible
membership in $b_1$ says nothing about the name's presence.
$\mathsf{Right}(f_2)$ optional: the code demands a $b_1$ index covering
the name with $\mathit{idx}_1.\mathsf{get} \leq f_2.\mathsf{get}$. If
$\mathsf{has}(o,\mathsf{name})$: the name is covered by
$\mathit{idx}_1$ and not named in $b_1$ (this is the $\mathsf{Right}$
case), so membership in $b_1$ puts the value in
$\sem{\mathit{idx}_1.\mathsf{get}} \subseteq \sem{f_2.\mathsf{get}}$;
if absent, the optional clause is vacuous. Rejecting when $b_1$ has no
covering index is again necessary: $o$ may carry the key with an
arbitrary value. For symbol fields the same two arms run with the symbol
index as the covering index --- same argument, one index kind.
\end{proof}

\subsection{Index pairs and the covering semantics}
\label{sec:indexes}

\begin{definition}[Index-Pair Check]
\texttt{isSub\_IndexPair}$(\mathit{idx}_1, \mathit{idx}_2)$:
$(-, \mathsf{None})$ accept; $(\mathsf{Some}, \mathsf{Some})$ require
get covariant and set contravariant;
$(\mathsf{None}, \mathsf{Some}(i_2))$ probe
$\top \leq i_2.\mathsf{get}$ and $i_2.\mathsf{set} \leq \bot$ --- which
fails for concrete index types and succeeds precisely when $i_2$'s
get/set are inference variables (the \emph{synthetic} index that object
literals carry), unifying them with $\top$/$\bot$.
The number-kind invocation uses \emph{effective} indexes: a side lacking
a number index falls back to its string index.
\impl{TTypeSubtyping.sk / isSub\_IndexPair, isSub\_ (number-index match);
TType.sk / TField::createSyntheticIndex, kSyntheticIndexName}
\end{definition}

\begin{lemma}[Index Covering Soundness]
\label{lem:indexes}
The three per-kind invocations establish the index clauses of
$o \in \sem{b_2}$, given the recursive inclusions.
\end{lemma}
\begin{proof}
Fix a prop $p \in \mathsf{props}(o)$ covered by an index of $b_2$ and
not named by $b_2$'s fields; we must place $\mathsf{get}(o,p)$ in that
index's get-denotation. If $p$ is named by $b_1$, the $\mathsf{Left}$
arm already did (Lemma~\ref{lem:width}). Otherwise $p$ is anonymous on
both sides and membership in $b_1$ constrains it only through $b_1$'s
indexes. Case symbol key: the symbol pair check gives
$\mathit{idx}^{\mathrm{sym}}_1.\mathsf{get} \leq
\mathit{idx}^{\mathrm{sym}}_2.\mathsf{get}$, and
$(\mathsf{None},\mathsf{Some})$ with concrete $i_2$ is rejected ---
necessary, as $o$'s symbol props are otherwise unconstrained; with
synthetic $i_2$, growth makes $\sem{i_2.\mathsf{get}} = \Val$ under the
accepted state, so the clause is vacuous and acceptance exact (the
sealed-to-$\bot$ set side makes the literal's index reject all writers
--- restrict-only). Case string-kind key: only $b_2$'s string index
covers $p$; the string pair check plus the same
$(\mathsf{None},\mathsf{Some})$ analysis concludes. Case numeric-kind
key: $p$ must satisfy \emph{every} covering index of $b_2$ --- its
number index and its string index. For $b_2$'s string index: $p$ is
covered on the $b_1$ side by $b_1$'s string index too (string covers
both kinds), and the string pair check transfers it. For $b_2$'s number
index: $o$'s constraint at $p$ from $b_1$ is its number index if
present, else its string index --- exactly the \emph{effective} number
index --- and the effective-pair check transfers it. The four-way match
in \texttt{isSub\_} enumerates precisely these fallbacks (both absent:
nothing to transfer, and if $b_2$ has only a string index the string
pair already handled numeric-kind keys). This is the point where the
main paper's move to the covering semantics pays off: under the older
per-kind-partition reading, the number-under-string fallback would have
been an unsound borrowing; under covering it is the exact transfer.
\end{proof}

\begin{corollary}[An Extensible Source Cannot Manufacture Closed Absence]
\label{cor:open-closed}
Subtyping cannot admit an object source that leaves an uncovered key
unconstrained into a \texttt{Closed<T>} target that proves the key absent.
This corollary establishes registry row \textbf{L-SUB-CLOSED-ABSENCE}
(conditional on the recursive \texttt{isSub} verdicts, as every field-stratum
result is).
\end{corollary}
\begin{proof}
\texttt{Closed<T>} represents every uncovered string or symbol key with a
concrete index whose $\mathsf{get}=\bot$ (number keys fall back to the string
index). Against a source with no corresponding index, the
$(\mathsf{None},\mathsf{Some}(i_2))$ arm must prove
$\top \leq i_2.\mathsf{get}=\bot$ and rejects. Against a source with a
concrete index, the $(\mathsf{Some},\mathsf{Some})$ arm requires
$i_1.\mathsf{get}\leq\bot$ and therefore accepts only when the source already
proves the same absence. A source may also \emph{name} an extra field the
closed target proves absent: the width lemma's extra-field case checks a
named source field with no target counterpart against the target's covering
index — here the never-valued one — so it must establish
$f_1.\mathsf{get}\leq\bot$ and rejects any inhabited extra field, string- and
symbol-named alike. Numeric-named extras consult the string never-index and
nothing else: \texttt{closeObjectType} rejects every \emph{declared} index
(string, number, or symbol; SKJS~9033) before installing the never-indexes,
and the synthetic number index is cleared — a valid closed target therefore
has no number index at all, so the single-index consultation of
L-SUB-IDXNAME cannot even arise here.
Closing also drops a \emph{synthetic} number index
(\texttt{withNeverIndexes}) — an object-literal inference artifact carrying
\texttt{TVar} get/set, describing no declared capability — and this is
absence-safe: with the number slot empty, numeric-kind keys fall under the
string never-index for \emph{admission} exactly as they do for the absence
claim, so the claim and the gate move together. Hence an accepted admission
never creates the absence fact on which the main paper's \textsc{R-In}
true-branch pruning relies. Note the structural contrast with
Finding~F-CABSENT: indexes have no lax middle state — an absent index reads
$\top$ and $\top \leq \bot$ rejects — so the two-step laundering that
defeats the callability brand has no analogue here. The pinned
open-to-closed control is
\texttt{tests/expect-error/closed-type.ts} /
\texttt{testOpenIntoClosedRejected}.
\impl{TEnv.sk / closeObjectType; TType.sk / withNeverIndexes,
createSyntheticIndex; TTypeSubtyping.sk / isSub\_IndexPair,
isSub\_NamedFields (extra-field arm)}
\end{proof}

\subsection{Callability: the lattice check and a violation}
\label{sec:modes}

\begin{definition}[Callability Check]
\label{def:callability-check}
\texttt{TFields::isSubCallability}$(m_1, m_2)$ over the
$\mathsf{Callability}$ lattice (\S\ref{sec:recap}), by arm
(\texttt{src/TTypeSubtyping.sk:1255-1297}):
$(-, \mathsf{CUnknown})$ accept;
$(\mathsf{CAbsent}, \mathsf{CAbsent})$ accept;
$(\mathsf{CUnknown}, \mathsf{CAbsent})$ accept --- \emph{lax by design},
analyzed below;
$(\mathsf{CPresent}, \mathsf{CAbsent})$ and
$(\mathsf{CFunction}, \mathsf{CAbsent})$ reject;
$(\mathsf{CPresent}, \mathsf{CFunction})$ and
$(\mathsf{CFunction}, \mathsf{CFunction})$ accept;
$(\mathsf{CAbsent}, \mathsf{CFunction})$,
$(\mathsf{CUnknown}, \mathsf{CFunction})$,
$(\mathsf{CFunction}, \mathsf{CPresent})$,
$(\mathsf{CAbsent}, \mathsf{CPresent})$, and
$(\mathsf{CUnknown}, \mathsf{CPresent})$ reject;
$(\mathsf{CPresent}(\sigma_1), \mathsf{CPresent}(\sigma_2))$ recurses
into \texttt{TFuns::isSub}. The same table serves \texttt{call} and
\texttt{new} (invocations at \texttt{:1240-1251}).
\impl{TTypeSubtyping.sk / TFields::isSubCallability}
\end{definition}

Against the semantics of \S\ref{sec:recap}, every arm except one is
exact or forced. The five rejections are semantically necessary
($\mathsf{callable}$ members cannot inhabit a
$\neg\mathsf{callable}$ claim, and nothing short of a proof of
$[[\mathrm{Call}]]$ may enter the positive $\mathsf{CFunction}$ brand
--- in particular $\mathsf{CUnknown} \not\leq \mathsf{CFunction}$, so
\emph{nothing launders into} $\mathsf{CFunction}$; likewise a bare
$\mathsf{CFunction}$ offers no signature to satisfy a
$\mathsf{CPresent}$ target). $(-, \mathsf{CUnknown})$ is exact: the
target claims nothing. $(\mathsf{CPresent}, \mathsf{CFunction})$ and
the $\mathsf{CFunction}$ diagonal are exact: the target claims
callability, which the source proves. The signature recursion is the
subject of \S\ref{sec:tfun}. The one lax arm is not sound:

\begin{finding}[F-CABSENT: The Laundered $\mathsf{CAbsent}$ Brand]
\label{find:cabsent}
$(\mathsf{CUnknown}, \mathsf{CAbsent})$ is accepted: an
unknown-callability source flows into a type whose mode claims
\emph{definitely not callable}. The arm exists for TS structural
assignability --- a plain object type flows into a class-instance type
(\texttt{\{a: number\}} into \texttt{C}), and class instances carry
$\mathsf{call} = \mathsf{new} = \mathsf{CAbsent}$
(\texttt{src/TEnv.sk:1003-1008}); the code comment names the cost
explicitly (\texttt{src/TTypeSubtyping.sk:1264-1271}). Against
\S\ref{sec:recap}'s semantics it is a genuine non-inclusion: a callable
$o \in \sem{b_1}$ (callables inhabit $\mathsf{CUnknown}$ branches ---
that is what makes the mode ``unknown'') is excluded from $\sem{b_2}$
by the $\mathsf{call}(b_2) = \mathsf{CAbsent} \Rightarrow
\neg\mathsf{callable}(o)$ clause. And the violated clause is
load-bearing in refinement: \texttt{narrowByTypeofFunction} drops
$\mathsf{CAbsent}$ branches in the \texttt{typeof x === "function"}
true case (\texttt{src/TType.sk:2438-2460}), so a callable laundered
\texttt{Function} $\to$ \texttt{\{\}} $\to$ \texttt{C} --- two
individually-accepted steps, $(\mathsf{CPresent}, \mathsf{CUnknown})$
then $(\mathsf{CUnknown}, \mathsf{CAbsent})$ --- is narrowed to
\texttt{never} in the branch the runtime takes. Witnessed end-to-end
(registry row \textbf{F-CABSENT}; issue draft
\path{issues/cabsent-brand-laundering.md}: SKJS accepts, Node takes the
branch with a function value at type \texttt{never}, tsc rejects ---
the unsoundness is SKJS-specific). $\mathsf{CAbsent}$ is thus a
\emph{brand}: a claim refinement trusts as proof but subtyping enforces
only against sources that themselves prove callability --- exactly the
trust model of the array brand (Finding~\ref{find:brand}), whose
callability instance this is. Repair options: (a) reject
$(\mathsf{CUnknown}, \mathsf{CAbsent})$ --- sound, breaks TS structural
assignability into class-instance types; (b) downgrade refinement's
trust: treat $\mathsf{CAbsent}$ as droppable under
\texttt{typeof "function"} only where every inhabitant is
class-constructed (no structural admission route) --- which is the same
decision the array brand needs, so one trust-model resolution covers
both findings. The same analysis applies verbatim to
\texttt{new}/\textsf{constructible}.
\impl{TTypeSubtyping.sk / isSubCallability (CUnknown, CAbsent) arm;
TType.sk / narrowByTypeofFunction; TEnv.sk / class-instance modes}
\end{finding}

\subsection{The array brand: an omission}

\begin{finding}[F-BRAND: \texttt{isArray} Is Never Compared]
\label{find:brand}
\texttt{TFields::isSub\_} contains no comparison of the
\texttt{isArray} fields. The brand is stripped out of the named-field
map at construction (\texttt{TFields::create} removes
\texttt{\_\_SKJS\_isArray} and folds it into the flag), so it is not
checked by the named-field arms either: a branch with
$\mathsf{isArr} = \mathsf{None}$ (or even $\mathsf{Some}(\mathsf{false})$)
is accepted against a branch with $\mathsf{isArr} =
\mathsf{Some}(\mathsf{true})$ whenever the structural checks pass.
Against the main paper's brand clause
($\mathsf{isArr}(b) = \mathsf{Some}(x) \Rightarrow \mathsf{isArrayV}(o)
= x$) this is a non-inclusion; and as with F-CABSENT the violated clause
is load-bearing elsewhere: \textsc{R-IsArray}'s \emph{false} branch drops
branches branded $\mathsf{Some}(\mathsf{true})$, so a structural
non-array impostor accepted at type
\texttt{ReadonlyArray<T>} is typed \texttt{never} in the
\texttt{else} of an \texttt{Array.isArray} test that its runtime takes.
Exploitability is narrower than F-CABSENT --- the impostor must
structurally satisfy every array member, which TS programs rarely do
accidentally, and tsc has no brand at all, so SKJS is here no
\emph{worse} than TS, merely less sound than its own model claims.
(The construction-time lowering of an array-branded branch's modes to
$\mathsf{CAbsent}$, \S\ref{sec:recap}, rests on this same brand: the
callability meet's $\bot$ under the array brand is exactly as
trustworthy as \texttt{isArray} itself.)
Repair: compare brands in \texttt{isSub\_}
($\mathsf{Some}(x) \leq \mathsf{Some}(x)$;
$\mathsf{None} \not\leq \mathsf{Some}(x)$), or downgrade the brand
clause to advisory and re-prove \textsc{R-IsArray} without the
$\mathsf{Some}(\mathsf{true})$-drop.
\impl{TTypeSubtyping.sk / TFields::isSub\_ (no isArray arm);
TType.sk / TFields::create (brand stripping)}
\end{finding}

\begin{theorem}[Field Stratum, Assembled]
\label{thm:tfields}
\texttt{TFields::isSub} is one-step sound for the main paper's object
membership \emph{with the $\mathsf{CAbsent}$ clause restricted to
callability-proving sources and the brand clause deleted},
in the read direction conditional only on \textbf{L-SUB-IDXNAME} at
the numeric-named-field/dual-index corner
(Lemmas~\ref{lem:both-get}, \ref{lem:width}, \ref{lem:indexes} — the
width lemma's named residual);
its set-direction checks are the \textbf{L-MUT} obligations
(Lemma~\ref{lem:both-set}); the two carve-outs are exactly the
F-CABSENT/F-BRAND violations, and the sound callability arms of
\S\ref{sec:modes} defer their signature recursion to
\S\ref{sec:tfun}. \qed
\end{theorem}

\section{Functions: \texttt{TFun::isSub} and a Behavioral Signature
Semantics}
\label{sec:tfun}

The main paper's membership semantics treats signatures as opaque, so
``\texttt{TFun::isSub} is sound'' is not even a proposition against it:
no set-inclusion consequence distinguishes right from wrong variance.
This section supplies the missing semantics --- the paper's main modeling
contribution --- proves the implementation's rules against it, and
proposes it as a named extension of the main paper's object semantics.

\begin{extension}[E-SIG: Call-Soundness Contract]
\label{ext:esig}
Let $\sigma = (\overline{\tau\ \mathrm{tparams}},\ \mathit{this}^?,\
\overline{p},\ \overline{q}^{\,\mathrm{opt}},\ r^{?\,\mathrm{rest}},\
\mathit{ret})$. Say an argument list $\vec{a}$ \emph{conforms to}
$\sigma$ if $|\vec{a}| \geq |\overline{p}|$ and, for every position $i$
that $\sigma$ \emph{covers} (the required $p_i$, then optional $q_j$,
then --- if a rest is present --- every further position at the rest
element type), $a_i \in \sem{\sigma_i}$; positions \emph{beyond}
$\sigma$'s coverage (possible only when $\sigma$ has no rest) are
unconstrained. Then define $o \vDash \sigma$ iff for every instantiation
$\theta$ of $\sigma$'s type parameters satisfying their constraints,
every $\mathit{this}$-value $t \in \sem{\theta\,\mathit{this}}$ (any $t$
if $\mathit{this}$ is absent), and every $\vec{a}$ conforming to
$\theta\sigma$: if calling $o$ with $(t, \vec{a})$ returns normally, the
result is in $\sem{\theta\,\mathit{ret}}$. The proposed refinement of
the main paper's object semantics: replace ``signatures are opaque'' by
--- $\mathsf{call}(b) = \mathsf{CPresent}(\sigma)$ entails
$\mathsf{callable}(o) \wedge o \vDash \sigma$;
$\mathsf{CFunction}$ entails $\mathsf{callable}(o)$ only (no signature
claim --- and \texttt{maybeGetCall} requires $\mathsf{CPresent}$, so the
checker never invokes through a bare $\mathsf{CFunction}$);
$\mathsf{CUnknown}$ entails nothing; dually for
$\mathsf{new}$ with construction. (The $\mathsf{CAbsent}$ clause is
unchanged from \S\ref{sec:recap} and its lax subtyping arm is graded as
Finding~\ref{find:cabsent}.) Two deliberate choices:
extra trailing arguments are \emph{unconstrained but permitted} ---
JS functions accept any arity, and the contract must quantify over them
or arity-widening subsumption would be unsound; and only normal returns
are constrained (exceptions untracked). In the stratified semantics of
\S\ref{sec:coinduction}, argument memberships on the left of the
implication are taken at full depth and the result membership at level
$n$, keeping the operator monotone.
\end{extension}

\begin{theorem}[Signature Variance is Sound for E-SIG]
\label{thm:tfun}
Suppose \texttt{TFun::isSub}$(f_1, f_2)$ accepts, with its recursive
pairs given (one-step). Then $o \vDash f_1 \Rightarrow o \vDash f_2$,
\emph{except} through the two arms isolated as
Findings~\ref{find:this} and~\ref{find:tparam} below.
\end{theorem}
\begin{proof}
Fix $o \vDash f_1$, an instantiation $\theta$ admissible for $f_2$, a
this-value $t$, and $\vec{a}$ conforming to $\theta f_2$; we must land
the result in $\sem{\theta\,\mathit{ret}_2}$.
\emph{Type parameters:} equal counts are enforced;
\texttt{RenameClauses} identifies $f_2$'s parameters with $f_1$'s, so
$\theta$ reads as an instantiation of the shared parameters; where both
constraints exist, the contravariant check
$\mathit{cstr}_2 \leq \mathit{cstr}_1$ (a recursive pair, transported by
Lemma~\ref{lem:clauses}'s parametricity through the renaming) makes
$\theta$ admissible for $f_1$; where neither exists, trivial; the
$(\mathsf{Some},\mathsf{None})$ arm is Finding~\ref{find:tparam}.
\emph{This:} the receiver is a four-state sum
(\texttt{ThisParam}: \texttt{NoThis}, \texttt{Explicit},
\texttt{ImplicitClass}, \texttt{ImplicitShape};
\texttt{src/TType.sk:1384-1411}; arms at
\texttt{src/TTypeSubtyping.sk:1992-2033}). A \texttt{NoThis} source
accepts outright --- its contract quantifies over every receiver. An
\texttt{ImplicitShape} source against a non-shape target rejects
(its receiver requirement is erased; nothing can be proven). The
generic arm is contravariant, $\mathit{this}_2 \leq \mathit{this}_1$,
with a \emph{missing} target receiver read as $\top$ --- a free-call
permit, so the source's declared \texttt{this} must accept anything;
a source demanding a real receiver is thereby rejected against a
receiver-less target (TS accepts that; SKJS rejects). This arm is
exactly E-SIG's quantification. The two remaining diagonals ---
\texttt{ImplicitShape}/\texttt{ImplicitShape} and
\texttt{ImplicitClass}/\texttt{ImplicitClass}, both accepted without
comparing receiver types --- are justified only under an
\emph{attachment} reading of the receiver that E-SIG does not express:
Finding~\ref{find:this}.
\emph{Arguments:} we claim $\vec{a}$ conforms to $\theta f_1$. Both
sides are first normalized by flattening tuple rests into positionals
(\texttt{resolveSpreadToTuple}; conformance is invariant under this
re-description of coverage, with the code's note that optionals
preceding a flattened tuple become positionally required).
$|\vec{a}| \geq |\overline{p}_2| \geq |\overline{p}_1|$, the latter
enforced. For each position $i$ covered by $f_1$: if $f_2$ covers $i$,
the pairing loops checked $\sigma_{2,i} \leq \sigma_{1,i}$
contravariantly --- required-vs-required/optional/rest, then
optional-vs-optional/rest, then rest-vs-rest and the
\texttt{STRecFunOptionalSpread} arm for $f_1$-optionals reached only by
$f_2$'s rest --- so $a_i \in \sem{\theta\sigma_{1,i}}$; if $f_2$ does
\emph{not} cover $i$, then $f_2$ has no rest and $i$ exceeds its
coverage, but precisely this case is closed off by the two tail checks:
$f_1$ may have neither more coverage
($|\overline{p}_1| + |\overline{q}_1| \leq
|\overline{p}_2| + |\overline{q}_2|$) nor a rest when $f_2$ lacks one
--- the two ``mysterious'' arity rejections exist exactly to keep
$f_1$-covered-but-$f_2$-unconstrained positions from arising. Positions
beyond $f_1$'s coverage are unconstrained by the contract. Hence the
call is within $o$'s contract for $\theta f_1$ and a normal result lies
in $\sem{\theta\,\mathit{ret}_1} \subseteq \sem{\theta\,\mathit{ret}_2}$
by the covariant return pair. The \texttt{cgNodes} edge recording on
success is a results-accumulator write, not a \texttt{SubstState} write
--- it does not perturb Part~I's audits.
\end{proof}

\begin{finding}[F-THIS-IMPLICIT: The Attach-Only Receiver Arms]
\label{find:this}
The \texttt{ImplicitShape}/\texttt{ImplicitShape} and
\texttt{ImplicitClass}/\texttt{ImplicitClass} diagonals accept without
comparing receiver types. Under E-SIG as stated they are unjustified:
the target's contract quantifies over receivers in
$\sem{\mathit{this}_2}$, and nothing relates $\mathit{this}_1$ to it
(for \texttt{ImplicitClass} the receivers may be outright incompatible
--- the arm exists so that generic-container variance and subclass
overrides survive, since a contravariant receiver check would reject
\texttt{ReadonlyArray<Dog>} $\leq$ \texttt{ReadonlyArray<Animal>}).
What justifies them is a different receiver model: an
\emph{attachment} semantics in which a method's receiver at a call is
the very object it was accessed on, so a method carried from a
$b_1$-view of an object to a $b_2$-view of the \emph{same} object still
runs on that object. The checker enforces the model's premise
separately --- shape methods are attach-only by construction (their
carrier type is erased and a detached use rejects), and call sites
check the receiver of \texttt{Explicit}/\texttt{ImplicitClass}
functions --- but E-SIG has no clause expressing ``receiver $=$
carrier'', so the two diagonals cannot be graded against it. Named
residual: extend E-SIG with an attachment clause and prove the
diagonals against it, or restrict them. Status recommendation:
\textsc{assumed} (the attachment premise), pinned tests recommended.
\impl{TTypeSubtyping.sk / TFun::isSub\_ (thisParam match);
TType.sk / ThisParam}
\end{finding}

\begin{finding}[F-TPARAM: Constraint Dropping]
\label{find:tparam}
$(\mathsf{Some}(\mathit{cstr}_1), \mathsf{None})$ on a type-parameter
pair is accepted: \texttt{<T extends string>(x: T) => T} passes as
\texttt{<T>(x: T) => T}. An $f_2$-caller may instantiate outside
$\mathit{cstr}_1$, exercising the sub's body off its contract; E-SIG
does not justify the arm. (The subsequent positional checks run with
the shared parameter \emph{rigid}, so they cannot recover the gap.)
Same severity shape as F-THIS. Status recommendation: named residual,
pinned test recommended.
\impl{TTypeSubtyping.sk / TFun::isSub\_ (tparam cstr match)}
\end{finding}

\noindent
With E-SIG ratified, the callability arms of \S\ref{sec:modes} other
than F-CABSENT's are also fully justified:
$(\mathsf{CPresent}(\sigma_1), \mathsf{CPresent}(\sigma_2))$ reduces to
$\sigma_1 \leq \sigma_2$ under Theorem~\ref{thm:tfun};
$(\mathsf{CPresent}, \mathsf{CFunction})$ and the
$\mathsf{CFunction}$ diagonal are exact, since the target claims
callability only and the source proves it (E-SIG assigns
$\mathsf{CFunction}$ no signature clause); and the rejections are
forced --- $\mathsf{CUnknown}$ does not entail callable, so nothing
unproven enters the positive brand, and a bare $\mathsf{CFunction}$
offers no signature to satisfy a $\mathsf{CPresent}$ target.

\section{Variance Plumbing: Why Each Side Expands the Way It Does}
\label{sec:variance}

Step (5) expands $u_1$ with \texttt{Covariant()} and $u_2$ with
\texttt{Contravariant()}; \texttt{expandToVectors} resolves each
\texttt{TVar} through \texttt{Subst::get} at that variance
(covariant: seal-and-return the binding, $\top$ if fresh; contravariant:
return the binding unsealed, $\bot$ if fresh), then applies the
variable's constraints via \texttt{TType::inter}.

\begin{lemma}[Directional Reads]
\label{lem:variance}
Validation is sound iff the sub-side expansion \emph{over}-approximates
and the super-side expansion \emph{under}-approximates: from
$\sem{u_1} \subseteq \sem{E_{\mathrm{cov}}(u_1)}$,
$\sem{E_{\mathrm{contra}}(u_2)} \subseteq \sem{u_2}$, and
$\sem{E_{\mathrm{cov}}(u_1)} \subseteq \sem{E_{\mathrm{contra}}(u_2)}$
(the validated inclusion) the target follows by composition; with either
direction flipped, the composition breaks --- e.g.\ reading a sub-side
open variable contravariantly would validate
$\alpha \leq \texttt{string}$ from a still-growing bound
$\texttt{"a"}$, later falsified when $\alpha$ grows to
$\texttt{"a"} \mid 1$; reading a fresh super-side variable covariantly
would seal it to $\top$ and accept everything against a variable whose
solution is not yet chosen. The implementation's assignment is the
correct one. \qed
\end{lemma}

The two approximation facts themselves decompose as follows.
\emph{Covariant/sub side:} Deep Dive~C's T-EXP proves expansion
over-approximates for modes \texttt{None}/\texttt{Covariant},
conditional on (among others) \textbf{L-INTER} and
\textbf{L-INTER-BOT} --- and Deep Dive~B
\emph{refuted} the $\bot$-clause \textbf{L-INTER-BOT} in its live strata
(incomparable signatures return $\bot$; the beyond-precision
num-spelling stratum). The
exposure here is concrete: \texttt{get}'s constraint application
$\mathsf{inter}(\mathit{binding}, \mathit{cstr})$ can under-shoot the true
intersection through those strata, shrinking the covariant expansion below
$\sem{\alpha}$ and validating vacuously. We name the residue:

\begin{obligation}[O-GETCOV: Covariant Reads Over-Approximate]
\label{obl:getcov}
$\sem{\alpha}_\rho \subseteq
\sem{\texttt{Subst::get}(\alpha, \mathsf{cov})}_{\rho'}$. Follows from
the instance-semantics axiom plus L-INTER/L-INTER-BOT; open on the remaining live
\textbf{L-INTER-BOT} strata (incomparable signatures;
beyond-precision num-spelling),
which the constraint application can still hit. Status: \textsc{open},
inherited.
\end{obligation}

\begin{obligation}[O-GETCONTRA: Contravariant Reads Under-Approximate]
\label{obl:getcontra}
$\sem{\texttt{Subst::get}(\alpha, \mathsf{contra})}_{\rho}
\subseteq \sem{\alpha}_{\rho}$, reading $\sem{\alpha}$ as the eventual
committed solution: $\bot$ for fresh variables is trivially below
anything; an $\SVOpen$ bound is below the final value because bounds
only grow (Deep Dive~E's transitions) --- \emph{but} the constraint
intersection applied on the way out must not \emph{enlarge} past
$\sem{\alpha}$. This is a meet lower-bound property, not the
over-approximation supplied by L-INTER/L-INTER-BOT, and the semantics of
``eventual solution'' for variables that are never sealed needs the main
paper to fix a convention. Deep Dive~C explicitly assigned this read to
L-SUB; we discharge its use in Lemma~\ref{lem:variance} and leave the
two flagged sub-claims open. Status: \textsc{open}.
\impl{TTypeSubst.sk / Subst::get; TType.sk / expandToVectors}
\end{obligation}

%==========================================================================
\section{Status and Recommended Ledger Edits}
\label{sec:status}

\subsection{Status by stratum}

\begingroup
\renewcommand{\arraystretch}{1.25}
\small
\begin{longtable}{@{}p{2.9cm}p{10.5cm}@{}}
\toprule
\textbf{Stratum / item} & \textbf{Result here} \\
\midrule
\endfirsthead
\toprule
\textbf{Stratum / item} & \textbf{Result here} \\
\midrule
\endhead
\bottomrule
\endlastfoot
Dispatch fast paths & Lemma~\ref{lem:fastpath} \\
Restore-on-failure & Lemma~\ref{lem:restore} (family-wide audit) \\
\texttt{subst.seenSub} & Theorem~\ref{thm:memo} (\S\ref{sec:substseen}) \\
\texttt{tenv.seenSub} & Lemmas~\ref{lem:version},
  \ref{lem:pathscope}, \ref{lem:coherence}: hit $=$ pending ancestor at
  identical bindings \\
Coinduction & Theorem~\ref{thm:coind}: L-SUB reduces to one-step strata \\
Head-depth refusal & Definition~\ref{def:headdepth}: rejection-only
  divergence guard, path-scoped \\
Growth pre-passes & guidance only (main paper theorem);
  MR-1 entry point noted; tuple positional seeding \\
Primitives, units & Lemma~\ref{lem:prim} \\
$\STypes$ covering & Lemma~\ref{lem:stypes} \\
$\Clauses$ DNF + var gate & Lemma~\ref{lem:clauses}, parametric in
  $\TParam$ \\
\texttt{TObject} branches & Lemma~\ref{lem:tobject} \\
Fields, read direction & Lemmas~\ref{lem:both-get}, \ref{lem:width},
  \ref{lem:indexes} (incl.\ covering indexes, synthetic-index probes,
  access/phantom/private) \\
Fields, write direction & Lemma~\ref{lem:both-set} \\
Callability lattice & Definition~\ref{def:callability-check}: all arms
  exact or forced except one \\
$(\mathsf{CUnknown},\mathsf{CAbsent})$ arm & Finding~\ref{find:cabsent}
  + executed witness \\
Array brand & Finding~\ref{find:brand} + composite witness \\
\texttt{TFun} variance & Theorem~\ref{thm:tfun} against E-SIG \\
\texttt{this}/tparam-cstr arms & Findings~\ref{find:this},
  \ref{find:tparam} \\
Variance plumbing & Lemma~\ref{lem:variance};
  Obligations~\ref{obl:getcov}, \ref{obl:getcontra} \\
Array-literal decision & Obligation~\ref{obl:arraylit} \\
\end{longtable}
\endgroup

\subsection{Stratum outcomes and recommendations}

\begin{enumerate}
\item \emph{The L-SUB surface decomposes as}: L-SUB-ARCH (caches, dispatch,
  coinduction; \textsc{proved} here mod A-HASH — except the
  \texttt{subst.seenSub} component, \textsc{violated} via the witnessed
  MR-1 until the memo is version-keyed); L-SUB-PRIM (\textsc{proved});
  L-SUB-DNF ($\STypes$/$\Clauses$/\texttt{TObject} skeletons;
  \textsc{proved} modulo lower strata); L-SUB-FIELDS-GET
  (\textsc{proved}); L-SUB-FIELDS-SET (conditional on L-MUT);
  L-SUB-MODES (\textsc{violated}, F-CABSENT); L-SUB-BRAND
  (\textsc{violated}, F-BRAND); L-SUB-FUN (\textsc{proved} against
  E-SIG, residuals F-THIS-IMPLICIT/F-TPARAM); L-SUB-EXPVAR
  (\textsc{open}, O-GETCOV/O-GETCONTRA); L-SUB-ARRAYLIT
  (\textsc{open}).
\item \emph{Recommended model extensions}:
  E-SIG (adopt Extension~\ref{ext:esig}, replacing ``signatures are
  opaque''); A-HASH (hash-consing identity) — both are registry
  entries. Note that adopting E-SIG strengthens the
  main paper's monotonicity axiom: its deflection of the contravariance
  counterexample (``membership never inspects parameter types'') stops
  being automatic and must cite the stratified-level bookkeeping of
  \S\ref{sec:coinduction}.
\item \emph{Resolve the brand-trust model explicitly}: the
  $\mathsf{CAbsent}$ and array brands are claims refinement trusts as
  proofs but subtyping enforces only against sources that themselves
  prove the property; either close the lax admissions
  ($(\mathsf{CUnknown},\mathsf{CAbsent})$; the missing brand arm) or
  downgrade the \textsc{R-Typeof}/\textsc{R-IsArray} drops to
  brand-insensitive narrowing --- one decision covers both findings;
  wire the executed F-CABSENT witness and F-BRAND's composite sketch as
  pinned tests either way.
\item \emph{Cross-reference}: the L-INTER-BOT row lists O-GETCOV as a
  downstream dependent, so that repairing the remaining L-INTER-BOT
  strata is understood to be on L-SUB's critical path, not only
  refinement's.
\end{enumerate}

\paragraph{Registry.}
The authoritative, current statuses are maintained in the registry
(\texttt{docs/ledger.yaml}), rendered as the main paper's Assumptions and
Proof Obligations tables; this paper is cited there as evidence. The
registry rows it supports:

% GENERATED by docs/tools/docs.py from docs/ledger.yaml — do not edit.

\begin{center}
\renewcommand{\arraystretch}{1.25}
\small
\begin{tabular}{@{}>{\raggedright\arraybackslash}p{2.4cm}>{\raggedright\arraybackslash}p{8.2cm}>{\raggedright\arraybackslash}p{2.6cm}@{}}
\toprule
\textbf{Row} & \textbf{Statement} & \textbf{Status} \\
\midrule
L-SUB & \texttt{Clauses::isSub} acceptance implies semantic inclusion under the current state. & \textsc{open} \\
E-SIG & Behavioral call-soundness semantics for function signatures (required for function subtyping to be a soundness statement). & \textsc{assumed} \\
A-HASH & \texttt{TType::compare} equality is trustworthy (XXHash-only fast paths). & \textsc{assumed} \\
L-SUB-FIELDS-SET & \texttt{isSub}'s field arms preserve write capability downward: acceptance implies each target field's $\mathsf{set}$-type is below the source's, and readonly fields acquire no writers. & \textsc{open} \\
L-SUB-CLOSED-ABSENCE & \texttt{isSub} admission cannot manufacture a \texttt{Closed<T>} absence fact: acceptance requires the source to prove $\mathsf{get}=\bot$ at every key the target proves absent. & \textsc{open} \\
O-GETCOV & Soundness of the covariant variable read in subtyping (inherits L-INTER's violated $\bot$-clause). & \textsc{open} \\
O-GETCONTRA & Soundness of the contravariant variable read in subtyping. & \textsc{open} \\
O-ARRAYLIT & Soundness of the array-literal decision procedure inside subtyping. & \textsc{open} \\
F-CABSENT & The $\mathsf{CAbsent}$ brand is trustworthy: refinement drops $\mathsf{CAbsent}$ branches under \texttt{typeof} \texttt{"function"}, but subtyping's lax $(\mathsf{CUnknown},\mathsf{CAbsent})$ arm admits callables into branded types. & \textsc{violated} \\
F-BRAND & Subtyping never compares the array brand, though refinement narrows on it. & \textsc{violated} \\
F-THIS-IMPLICIT & Subtyping between two shape receivers, or two class receivers, with implicit \texttt{this} parameters accepts without comparing their receiver types. & \textsc{assumed} \\
L-SUB-IDXNAME & A numeric-named field is checked against both a number and a string index when the supertype has both (soundness of the width lemma's numeric-named case). & \textsc{open} \\
A-INTERSEM & $\sem{\mathsf{TInter}(\bar t)} \subseteq \bigcap \sem{t}$ — the instance-layer upper bound the dispatcher's fast paths assume. & \textsc{assumed} \\
L-SUB-TSINST & The reflexive \texttt{TSInst} fast path in \texttt{isSub} (same name, arguments equal after \texttt{peekRaw}) is sound. & \textsc{violated} \\
\bottomrule
\end{tabular}
\end{center}


\paragraph{Imports (review contract).}
The named statements this paper assumes from elsewhere. A reviewer treats
each as an axiom, checking citations against the registry's canonical
statement; no other paper's internals are in scope:

% GENERATED by docs/tools/docs.py from docs/ledger.yaml — do not edit.

\begin{center}
\renewcommand{\arraystretch}{1.25}
\small
\begin{tabular}{@{}>{\raggedright\arraybackslash}p{2.6cm}>{\raggedright\arraybackslash}p{7.4cm}>{\raggedright\arraybackslash}p{1.2cm}>{\raggedright\arraybackslash}p{2.2cm}@{}}
\toprule
\textbf{Row} & \textbf{Canonical statement} & \textbf{Owner} & \textbf{Status} \\
\midrule
L-INTER & \texttt{inter} over-approximates set intersection (the containment clause). & B & \textsc{proved} \\
L-INTER-BOT & A $\bot$ result from \texttt{inter} implies semantic disjointness. & B & \textsc{violated} \\
L-EXP & Expansion preserves or enlarges semantics (modes \texttt{None}/Covariant); residual instances keep theirs. & C & \textsc{proved} \\
A-CTOR-CALL & Class-constructor values carry $\mathsf{CPresent}$ construct and $\mathsf{CFunction}$ call modes, consistent with $[[\mathrm{Call}]]$/$[[\mathrm{Construct}]]$. & A & \textsc{proved} \\
L-LIT & Literal spellings: (i) $\delta$ is a function; (ii) injectivity modulo $\approx$ after normalization; (iii) falsiness-exactness of the lexeme test. & G & \textsc{violated} \\
L-MUT & Field writes are governed by $\mathsf{set}$-types; get-only membership is sound for snapshot reads. & main & \textsc{assumed} \\
OBL-MONOTONE & Monotone Reads Under Growth: judgements validated under $\rho$ survive later legal transitions. & E & \textsc{violated} \\
OBL-RESTORE & No Commitment on Failure: failing and exploratory paths restore all relevant state. & E & \textsc{proved} \\
MR-1 & No \texttt{subst.seenSub} memo recorded under \texttt{freezeCandidates} may outlive growth of an open candidate readable from its left side. & E & \textsc{violated} \\
E-ARR & Delayed array-literal forms keep their semantics under expansion. & C & \textsc{assumed} \\
SEM-GFP & The set-theoretic semantics is the greatest solution of the recursive membership equations. & main & \textsc{proved} \\
E-TRANSPORT & Monotone-reads transport: an inclusion validated under $\rho$ survives later legal transitions on the same timeline, except for the cache behavior recorded by MR-1. & E & \textsc{proved} \\
\bottomrule
\end{tabular}
\end{center}


\subsection{Failed-attempt log}

Recorded per series policy: (i) the global ``version-as-state-hash''
argument for \texttt{tenv.seenSub} fails (two concrete counterexample
schemas, \S\ref{sec:naive}); the successful replacement needs the frozen
\texttt{TEnv} threading and save/restore locality, neither of which is
about the version at all. (ii) A membership-only proof of the field
stratum fails at exactly two clauses (the $\mathsf{CAbsent}$ clause
under its lax admission arm, and the brand) --- both times
because the main paper priced those clauses into \emph{refinement}
proofs; the failures are reported as violations rather than weakened
into assumptions. (iii) A proof of \texttt{TFun::isSub} against the
opaque-signature semantics is not merely hard but type-incorrect --- the
conclusion is not expressible; hence E-SIG. (iv) An unconditional proof
of the covariant read (O-GETCOV) fails against Deep Dive~B's L-INTER-BOT
counterexamples; we did not find a validator-local workaround, since the
constraint application lives inside \texttt{Subst::get}.

\end{document}
